\chapter{جداول الصفحات}
\label{CH:MEM}

جداول الصفحات هي الميكانيكية الأساسية التي يقدمها العتاد وتعتمد عليها النواة لتوفير الذاكرة الافتراضية وتحقيق العزل بين العمليات.
تسمح جداول الصفحات لنظام التشغيل برسم خريطة العناوين الافتراضية المشتغلة عليها البرامج وتحويلها إلى عناوين فيزيائية فعلية في ذاكرة الحاسوب (RAM).

يقدم هذا الفصل كيفية تنظيم جداول الصفحات في معمارية RISC-V Sv39، وتخطيط الذاكرة لكل من النواة والعمليات في \texttt{xv6}، وكيفية تخصيص وتحرير الذاكرة الفيزيائية.

\section{عتاد التصفح}
\label{sec:paging_hardware}

تستخدم معمارية \textsf{RISC-V 64-bit} نمط تصفح يُدعى \textbf{Sv39}، حيث تُستخدم 39 بتًا فقط من أصل 64 بت في العنوان الافتراضي.
يتم تقسيم الذاكرة إلى كتل ثابتة تُسمى \textbf{الصفحات} (Pages)، ويبلغ حجم الصفحة الواحدة 4096 بايت (4 كيلوبايت).

يتكون العنوان الافتراضي ذو الـ 39 بت في Sv39 من:
\begin{itemize}
  \item 27 بت تعبر عن \textbf{رقم الصفحة الافتراضية} (Virtual Page Number - VPN)، وتُقسم إلى ثلاثة مستويات بعرض 9 بتات لكل مستوى (\texttt{VPN[2]}، \texttt{VPN[1]}، \texttt{VPN[0]}).
  \item 12 بت تعبر عن \textbf{الإزاحة} (Offset) داخل الصفحة ($2^{12} = 4096$ بايت).
\end{itemize}

يتكون جدول الصفحات في Sv39 من هيكل شجري ثلاثي المستويات:
\begin{enumerate}
  \item يحتفظ السجل العتادي \texttt{satp} (Supervisor Address Translation and Paging) بالعنوان الفيزيائي لجدول الصفحات من المستوى الأعلى (Level 2).
  \item تستخدم وحدة إدارة الذاكرة (MMU) الـ 9 بتات الأهم (\texttt{VPN[2]}) لفهرسة الجدول الأول للحصول على عنوان جدول المستوى التالي.
  \item تُستخدم الـ 9 بتات الوسطى (\texttt{VPN[1]}) لفهرسة جدول المستوى الثاني.
  \item تُستخدم الـ 9 بتات الأدنى (\texttt{VPN[0]}) لفهرسة جدول المستوى الثالث للحصول على \textbf{مدخل جدول الصفحات} (Page Table Entry - PTE).
\end{enumerate}

يحتوي مدخل جدول الصفحات (PTE) ذو الـ 64 بت على رقم الصفحة الفيزيائية (Physical Page Number - PPN) ورايات التحكم والتحقق (Flags):
\begin{itemize}
  \item \texttt{PTE\_V} (Valid): يحدد ما إذا كان المدخل صالحًا وموجودًا.
  \item \texttt{PTE\_R} (Readable): يسمح بقراءة الصفحة.
  \item \texttt{PTE\_W} (Writable): يسمح بالكتابة في الصفحة.
  \item \texttt{PTE\_X} (Executable): يسمح بتنفيذ التعليمات من الصفحة.
  \item \texttt{PTE\_U} (User): يحدد إمكانية الوصول للصفحة من وضع المستخدم (U-mode).
\end{itemize}

لتسريع ترجمة العناوين وتجنب الوصول المكرر للذاكرة، يحتوي المعالج على «مخزن الترجمة المؤقت (Translation Lookaside Buffer، TLB)». وعند تعديل جدول الصفحات، يجب على النواة تفريغ مخزن الترجمة المؤقت باستخدام تعليمة \texttt{sfence.vma}.

\section{فضاء عناوين النواة}
\label{sec:kernel_address_space}

يقوم \texttt{xv6} بإعداد \textbf{فضاء عناوين النواة} (Kernel Address Space) لتمكين النواة من الوصول إلى العتاد والأجهزة الملحقة والذاكرة المادية.
يستخدم \texttt{xv6} التعيين المباشر (Direct Mapping) لمعظم الذاكرة الفيزيائية داخل النواة:
\begin{itemize}
  \item العناوين من \texttt{0x00000000} إلى \texttt{0x80000000} تُخصص لمتحكمات الأجهزة والملحقات العتادية (مثل UART، PLIC، ومتحكمات virtio).
  \item العنوان \texttt{0x80000000} هو بداية الذاكرة الفيزيائية الرئيسية (RAM). تُعين النواة كودها وبياناتها بدءًا من هذا العنوان.
  \item يتواجد أقصى عنوان في الذاكرة المادية عند \texttt{PHYSTOP} (\texttt{0x88000000}، بحد أقصى 128 ميجابايت في \texttt{xv6}).
\end{itemize}

هناك استثناءان مهمان للتعيين المباشر في فضاء عناوين النواة:
\begin{enumerate}
  \item \textbf{صفحة المنصة القافزة} (Trampoline Page): تُعين عند أعلى عنوان افتراضي ممكن (\texttt{MAXVA - PGSIZE}) لتمكين التحويل بين فضاء المستخدم وفضاء النواة.
  \item \textbf{مكدسات النواة} (Kernel Stacks): تُخصص لكل عملية مكدس نواتي خاص مع وضع \textbf{صفحة حارس} (Guard Page) غير صالحة تحته مباشرةً لاكتشاف طفح المكدس (Stack Overflow) والتسبب في خطأ حماية فورية.
\end{enumerate}

\section{الكود البرمجي: إنشاء فضاء العناوين}
\label{sec:code_creating_address_space}

تُدار جداول الصفحات في \texttt{xv6} عبر البرمجية \texttt{kernel/vm.c}.
الدالة الأساسية للبحث والتنقل في شجرة جداول الصفحات هي \texttt{walk()}:
\begin{lstlisting}
pte_t *
walk(pagetable_t pagetable, uint64 va, int alloc)
\end{lstlisting}
تقوم \texttt{walk()} بتطبيق عملية الترجمة ثلاثية المستويات عتاديًا وبرمجيًا. إذا كان خيار \texttt{alloc} مفعلًا ولم يكن جدول الصفحات الفرعي موجودًا، تقوم \texttt{walk()} بتخصيص صفحة جديدة باستخدام \texttt{kalloc()}.

وتوفر الدالة \texttt{mappages()} آلية ربط نطاق من العناوين الافتراضية بالعناوين الفيزيائية عن طريق إضافة المداخل المناسبة عبر \texttt{walk()}.

\section{تخصيص الذاكرة الفيزيائية}
\label{sec:physical_memory_allocation}

تحتاج النواة إلى تخصيص وتحرير صفحات الذاكرة المادية بكتل حجمها 4096 بايت أثناء التشغيل لخدمة جداول الصفحات ومكدسات العمليات والذاكرات المؤقتة.
تُدار الذاكرة الفيزيائية الحرة بين نهاية كود النواة (\texttt{end}) وأقصى الذاكرة (\texttt{PHYSTOP}) بواسطة \textbf{مخصص الذاكرة الفيزيائية} (Physical Memory Allocator).

\section{الكود البرمجي: مخصص الذاكرة الفيزيائية}
\label{sec:code_physical_memory_allocator}

يتواجد مخصص الذاكرة في \texttt{kernel/kalloc.c}.
يحتفظ المخصص بقائمة متصلة أحادية من الصفحات الحرة (\texttt{struct run}):
\begin{lstlisting}
struct run {
  struct run *next;
};

struct {
  struct spinlock lock;
  struct run *freelist;
} kmem;
\end{lstlisting}
تضمن \texttt{spinlock} التزامن ومنع حالة التسابق (Race Condition) عند تخصيص أو تحرير الصفحات بين أنوية المعالج المختلفة:
\begin{itemize}
  \item \texttt{kalloc()}: تسحب صفحة حرّة من رأس القائمة \texttt{freelist} وتعيد عنوانها الفيزيائي بعد تعبئته بقيمة مخفية خردة (\texttt{1}) لاكتشاف الأخطاء.
  \item \texttt{kfree(pa)}: تتأكد من صحة العنوان الفيزيائي، وتكتب فيه بيانات خردة (\texttt{1}) لمنع استخدام الذاكرة المحررة، ثم تعيد الصفحات إلى أصل القائمة \texttt{freelist}.
\end{itemize}

\section{فضاء عناوين العملية}
\label{sec:process_address_space}

تملك كل عملية في \texttt{xv6} جدول صفحات مستقل يمثل فضاء عناوينها الافتراضي.
يبدأ فضاء عناوين العملية من العنوان الافتراضي \texttt{0x0}:
\begin{itemize}
  \item التعليمات والبيانات والمكدس الابتدائي تبدأ من العنوان \texttt{0x0}.
  \item ينظر مكدس المستخدم إلى الأسفل ويتم توسيعه تلقائيًا.
  \item تعلو القمة المتحركة للذاكرة الديناميكية للعملية (\texttt{heap}) وتُوسع بواسطة استدعاء النظام \texttt{sbrk()} المعتمد على الدالة \texttt{growproc()}.
  \item في أعلى فضاء العناوين الممكن (\texttt{MAXVA})، تقع صفحة المنصة القافزة (\texttt{TRAMPOLINE}) وتحتها مباشرة صفحة إطار المصيدة (\texttt{TRAPFRAME}).
\end{itemize}

\section{الكود البرمجي: exec}
\label{sec:code_exec}

يقوم استدعاء النظام \texttt{exec()} بالخطوات التالية لتشغيل برنامج جديد:
\begin{enumerate}
  \item يفتح الملف التنفيذي المنسق بصيغة ELF باستخدام \texttt{namei()}.
  \item يتفحص رأس ELF للتأكد من الماجيك هيدر الصحيح.
  \item ينشئ جدول صفحات جديد فارغ للعملية باستخدام \texttt{proc\_pagetable()}.
  \item يحمل أجزاء البرنامج (Program Segments) إلى الصفحات الفيزيائية المخصصة حديثًا عبر \texttt{loadseg()}.
  \item يخصص صفحة مكدس المستخدم ويدفع فيها الوسائط (\texttt{argv}).
  \item يحرر جدول الصفحات القديم والذاكرة القديمة عبر \texttt{proc\_freepagetable()}.
  \item يربط جدول الصفحات الجديد بالعملية ويثبت نقطة الدخول \texttt{epc}.
\end{enumerate}

\section{العالم الحقيقي}
\label{sec:real_world_mem}

تستخدم أنظمة التشغيل الحديثة آليات ذاكرة افتراضية متطورة للغاية مثل:
\begin{itemize}
  \item \textbf{النسخ عند الكتابة} (Copy-on-Write - COW): مشاركة الصفحات الفيزيائية بين الأب والابن عند \texttt{fork()} ونساخها فقط عند محاولة التعديل.
  \item \textbf{الصفحات الضخمة} (Superpages): استخدام صفحات بحجم 2 ميجابايت أو 1 جيجابايت لتقليل الضغط على مخازن الترجمة المؤقتة.
  \item \textbf{التصفح بالمطالبة} (Demand Paging) والتبديل (Swapping): نقل الصفحات إلى القرص الصلب عند امتلاء الذاكرة الفيزيائية.
\end{itemize}

\section{تمارين}
\label{sec:exercises_mem}

\begin{enumerate}
  \item اشرح لماذا تحتاج النواة إلى صفحة المنصة القافزة \texttt{TRAMPOLINE} في نفس العنوان الافتراضي لدى جميع العمليات والنواة.
  \item قم بتنفيذ آليات تخصيص الصفحات الكبيرة لتقليل عدد مستويات جداول الصفحات في \texttt{xv6}.
\end{enumerate}
